\newpage
\section{第17章：子类型化的 ML 实现}

\begin{taplauthorsolutionentrybox}
\phantomsection
\label{sol:author:17.3.1}
\textbf{17.3.1 解答}

只需把\cref{prop:16.3.2}中的算法写成程序：

\begin{taplocamlcode}
let rec join tyS tyT =
  match (tyS,tyT) with
    (TyArr(tyS1,tyS2),TyArr(tyT1,tyT2)) ->
      (try TyArr(meet tyS1 tyT1, join tyS2 tyT2)
       with Not_found -> TyTop)
  | (TyBool,TyBool) ->
      TyBool
  | (TyRecord(fS), TyRecord(fT)) ->
      let labelsT = List.map (fun (li,_) -> li) fT in
      let commonLabels =
        List.filter (fun l -> List.mem l labelsT)
          (List.map (fun (li,_) -> li) fS) in
      TyRecord
        (List.map
          (fun li ->
            (li, join (List.assoc li fS) (List.assoc li fT)))
          commonLabels)
  | _ ->
      TyTop
\end{taplocamlcode}
\taplrustcounterpart{sol-author-17-join}

\begin{taplocamlcode}
and meet tyS tyT =
  match (tyS,tyT) with
    (TyTop,tyT) -> tyT
  | (tyS,TyTop) -> tyS
  | (TyArr(tyS1,tyS2),TyArr(tyT1,tyT2)) ->
      TyArr(join tyS1 tyT1, meet tyS2 tyT2)
  | (TyBool,TyBool) ->
      TyBool
  | (TyRecord(fS), TyRecord(fT)) ->
      let labelsS = List.map (fun (li,_) -> li) fS in
      let labelsT = List.map (fun (li,_) -> li) fT in
      let allLabels =
        labelsS @ List.filter (fun l -> not (List.mem l labelsS)) labelsT in
      TyRecord
        (List.map
          (fun li ->
            if List.mem li labelsS && List.mem li labelsT then
              (li, meet (List.assoc li fS) (List.assoc li fT))
            else if List.mem li labelsS then
              (li, List.assoc li fS)
            else
              (li, List.assoc li fT))
          allLabels)
  | _ ->
      raise Not_found
\end{taplocamlcode}
\taplrustcounterpart{sol-author-17-meet}

再为条件式补上类型检查分支：

\begin{taplocamlcode}
| TmIf(fi,t1,t2,t3) ->
    if subtype (typeof ctx t1) TyBool then
      join (typeof ctx t2) (typeof ctx t3)
    else
      error fi "guard of conditional not a boolean"
\end{taplocamlcode}
\taplrustcounterpart{sol-author-17-conditional-type}

\taplsolutionbacklink{ex:17.3.1}{返回练习17.3.1}
\end{taplauthorsolutionentrybox}

\begin{taplauthorsolutionentrybox}
\phantomsection
\label{sol:author:17.3.2}
\textbf{17.3.2 解答}

见原书配套源码中的 \texttt{rcdsubbot} 实现。它在类型语法中加入
\texttt{TyBot}，在子类型检查器中令 \texttt{TyBot} 成为每种类型的子类型，
并按照\taplsecref{16.4}{sec:16.4}为应用和投影加入函数项或记录项为
\texttt{TyBot} 时的分支。

\taplsolutionbacklink{ex:17.3.2}{返回练习17.3.2}
\end{taplauthorsolutionentrybox}
